% Cartan-data obstruction for CHL levels N=4 and N=6.
% Author: Raeez Lorgat.

\providecommand{\kBKM}{\kappa_{\mathrm{BKM}}}
\providecommand{\kch}{\kappa_{\mathrm{ch}}}
\providecommand{\kcat}{\kappa_{\mathrm{cat}}}
\providecommand{\kfib}{\kappa_{\mathrm{fiber}}}
\providecommand{\fgDel}{\mathfrak{g}_{\Delta_5}}
\providecommand{\ZZ}{\mathbb{Z}}
\providecommand{\QQ}{\mathbb{Q}}
\providecommand{\LIIN}{\Lambda^{2,1}_{\mathrm{II},N}}

\section*{Cartan data and CHL levels $N=4,6$}
\label{sec:cartan-N46-obstruction}

\paragraph{Borcherds denominator constants.}
For the CHL paramodular ladder
\[
  N \in \{1,2,3,4,6\},
\]
the Borcherds-input constants and BKM modular characteristics are
\[
  (c_1(0),c_2(0),c_3(0),c_4(0),c_6(0))
  = (10,8,6,4,2),
\]
\[
  \bigl(\kBKM(\Phi_1),\kBKM(\Phi_2),\kBKM(\Phi_3),
  \kBKM(\Phi_4),\kBKM(\Phi_6)\bigr)
  = \frac{1}{2}(10,8,6,4,2)
  = (5,4,3,2,1).
\]
Thus $c_4(0)=4$ and $\kBKM(\Phi_4)=2$, while $c_6(0)=2$ and
$\kBKM(\Phi_6)=1$. These are Borcherds-weight data. A higher-level
rank-three BKM presentation requires further data: a real-simple-root
system, a Weyl vector, a parity splitting, and a denominator identity.

The constructed rank-three BKM superalgebra is $\fgDel$, attached to
the level-one form $\Delta_5=\Phi_1$.
For $N=4,6$, the notation $\mathfrak{g}^{(N)}$ names a target only
after the following data are supplied in one normalisation: a hyperbolic
root lattice, real and imaginary simple roots, a Weyl chamber, a Weyl
vector, parity of the imaginary roots, multiplicities from the Fourier
coefficients, and the Borcherds denominator identity. A fixed-subalgebra
formula $\mathfrak{g}^{(N)}=\fgDel^{\ZZ/N}$ leaves these data
unspecified.

\paragraph{Determinants and weights.}
The formula
\[
  \det A^{(N)} = 2(c_N(0)-2)
\]
mixes two different invariants. With the CHL constants above it gives
\[
  2(c_4(0)-2)=4,\qquad 2(c_6(0)-2)=0,
\]
whereas the displayed symmetric Gram matrices
\[
  G^{(4)} =
  \begin{pmatrix}
    2 & -1 & -2\\
    -1 & 2 & -2\\
    -2 & -2 & 4
  \end{pmatrix},
  \qquad
  G^{(6)} =
  \begin{pmatrix}
    2 & -1 & -2\\
    -1 & 2 & -2\\
    -2 & -2 & 6
  \end{pmatrix}
\]
have
\[
  \det G^{(4)}=-12,\qquad \det G^{(6)}=-6,
\]
and signature $(2,1)$ in both cases. The determinants are lattice
discriminants of these displayed forms, not consequences of
$\kBKM(\Phi_N)=c_N(0)/2$.

\paragraph{Cartan-integrality test.}
If $G^{(N)}$ is read as the Gram form of real simple roots
$\alpha_1,\alpha_2,\alpha_3$, the generalized Cartan entries are
\[
  a^{(N)}_{ij}
  = \frac{2(\alpha_i,\alpha_j)}{(\alpha_i,\alpha_i)}
  = \frac{2G^{(N)}_{ij}}{G^{(N)}_{ii}} .
\]
This gives
\[
  A^{(4)} =
  \begin{pmatrix}
    2 & -1 & -2\\
    -1 & 2 & -2\\
    -1 & -1 & 2
  \end{pmatrix},
  \qquad
  A^{(6)} =
  \begin{pmatrix}
    2 & -1 & -2\\
    -1 & 2 & -2\\
    -\frac{2}{3} & -\frac{2}{3} & 2
  \end{pmatrix}.
\]
The $N=6$ matrix has nonintegral off-diagonal Cartan entries and lies
outside the standard real-simple-root definition of a Borcherds
generalized Cartan matrix. The $N=4$ matrix passes this elementary
integrality test, which supplies an integral symmetrizable Cartan matrix
only. Imaginary simple roots, parity, multiplicities, and a denominator
identity remain additional construction data for both levels.

\paragraph{The coordinate root model is discarded.}
The coordinate formula
\[
  \delta_1=f_2+f_{-2}-f_3,\qquad
  \delta_2=f_{-2}+f_2-f_3
\]
has $\delta_1=\delta_2$. It cannot span a rank-three lattice and cannot
realise either $G^{(4)}$ or $G^{(6)}$. The accompanying expression for
$(\delta_3,\delta_3)$ also mixes incompatible scaling terms. No
coordinate embedding of the proposed roots into $\LIIN$ is part of the
proved statement.

\paragraph{Formal Weyl-vector calculation.}
Solving only the linear equations
\[
  G^{(N)}\rho^{(N)}
  = -\frac{1}{2}\operatorname{diag}(G^{(N)})
\]
gives
\[
  \rho^{(4)}
  = 2\alpha_1+2\alpha_2+\frac{3}{2}\alpha_3,\qquad
  (\rho^{(4)},\rho^{(4)})=-7,
\]
\[
  \rho^{(6)}
  = 6\alpha_1+6\alpha_2+\frac{7}{2}\alpha_3,\qquad
  (\rho^{(6)},\rho^{(6)})=-\frac{45}{2}.
\]
These are formal solutions for the displayed symmetric forms. They
become Weyl vectors for BKM denominator formulae only after the root
lattice, positive cone, imaginary simple roots, root multiplicities,
and parity have been fixed.

\paragraph{Level-six geometric anchors.}
At $N=6$, the CHL class uses the Niemeier anchor $6D_4$. The cyclic
quotient surface $K3/\langle g_6\rangle$ has local Du Val type
\[
  2A_5+2A_2+2A_1,
\]
obtained from the fixed-point counts
\[
  \#\operatorname{Fix}(g_6)=2,\qquad
  \#\operatorname{Fix}(g_6^2)=6,\qquad
  \#\operatorname{Fix}(g_6^3)=8.
\]
The Euler check is
\[
  e(K3/\ZZ_6)=\frac{24+2+6+8+6+2}{6}=8,\qquad
  8+(2\cdot 5+2\cdot 2+2\cdot 1)=24.
\]
The Niemeier lattice $6D_4$ is a global positive-definite even
unimodular rank-$24$ lattice. The singularity type
$2A_5+2A_2+2A_1$ is a local analytic stratification of the quotient
surface. The lattice $4A_5+D_4$ is another Niemeier lattice of Coxeter
number $6$. The $6A$ CHL anchor is $6D_4$.

\paragraph{Invariant separation.}
The quantities
\[
  \kBKM,\qquad \kch,\qquad \kcat,\qquad \kfib
\]
belong to different constructions. The constants above compute only
$\kBKM(\Phi_N)$, the weight of the chosen Borcherds product. The
$\kch$-, K\"unneth-multiplicative $\kcat$-, and Mukai-rank $\kfib$-
values are separate invariants. For $d\geq 3$ the native output of
$\Phi_d$ is $E_1$-chiral; any CHL/Borcherds denominator comparison is a
further specialisation with its own automorphic input.

\paragraph{Result.}
The proved statement is the Borcherds-weight identity
\[
  \kBKM(\Phi_N)=\frac{c_N(0)}{2}
  \quad (N\in\{1,2,3,4,6\}),
\]
with $(c_N(0))=(10,8,6,4,2)$. A rank-three BKM presentation attached to
the displayed $N=4,6$ matrices requires data absent from this note: an
integral Cartan normalisation at $N=6$, distinct coordinate roots,
imaginary roots, parity, multiplicities, and a denominator identity.

\paragraph{References.}
Borcherds 1998, \emph{Inventiones Mathematicae} 132, Theorem 13.3;
Gritsenko--Nikulin 1995, Part II, Theorem 2.1; Gritsenko 1999,
\emph{Arithmetical lifting and its applications}; Mukai 1988,
\emph{Inventiones Mathematicae} 94, Table 2; Nikulin 1979, Theorem 4.5;
Cheng--Duncan--Harvey 2014, umbral Niemeier tables; the Vol III
universal Borcherds-weight theorem, closed-form $c_N(0)$ lemma, and
Niemeier-versus-fixed-locus remark in the $K3\times E$ chapter.
